% Transcription — fonds Alexandre Grothendieck, shelfmark 16, batch 2 (pages 21-40).
% Established from the facsimile archives/batches/16/batch-02.pdf (a mirror of
% https://grothendieck.umontpellier.fr/16.pdf), per the skill
% transcribe-grothendieck. The batch file carries no cover sheet: PDF page N is
% archivists' page N + 20 (the pencilled numbers 21, 32, 33, 35, 37, 39 confirm
% it where they are legible).
%
% Pass: Opus 5 (claude-opus-5), 2026-09-19 — first pass, unchecked against the pages by a human.
%
% What the batch is. Nine leaves carry mathematics; the rest are unrelated
% typescripts used as paper (pages 22, 25, 27, 29, 34, 36, 38, 40 — the Brauer
% group, corrections to SGAD, a bibliography, English pages on group schemes
% and on injective sheaves, Karoubi on Banach categories) and one blank
% (page 23). They are skipped without further mention.
%
% Three runs.
%
%   21      A fast pencil draft, most of it cancelled: the computation
%           L^n_0 = L^n_0 π_0 = L^n_0 Λ'^n L'^n = αβ L^n_0, whence αβ = 1; a
%           remark on what the Lefschetz condition gives for L^n and Λ̄; then
%           an algebra of dimension 6 in degrees 0 to 4, with relations
%           L^4 = 0, L'^2 − LL' = 0, L^2L' − L^3 = 0, L^3L' = 0, where L is
%           Lefschetz and L' is not although L'^3 = L^3. The connective prose
%           is largely lost; the formulas carry.
%   24-28   « Algèbres de Hodge » — the title is his, taken from the otherwise
%           blank leaf page 30 (that leaf gets no page number of its own).
%           Page 24 fixes a graded algebra H^0…H^4 = k, V, kξ+W, V̌, kξ^2 and
%           the data φ, ψ, Q that determine its multiplication; the
%           associativity condition, the Lefschetz condition (φ non
%           dégénérée) and the Poincaré condition (Q non dégénérée). Page 26
%           treats the case (xy)(zt) = 0 and then recasts the whole datum as
%           (V, χ, (W̃, Q̃, α), ξ). Page 28 is the example χ = φ^2 with α = id,
%           and the orbit of ξ and φ̌ under an endomorphism u of Λ^2 V.
%   32-39   « Formulaire de L, Λ » — his title again, from the blank leaf
%           page 31 and from the typescript's own heading. Page 32 is a
%           typescript, heavily corrected in his hand (whole blocks cancelled
%           and replaced, index ranges overwritten in ink, the operator T_X
%           introduced by (6)); pages 33 to 39 are the handwritten
%           continuation, which proves the corollary (A(X×X) et C(Y)) implique
%           C(X) and, on page 39, the converse direction C(X) ⟹ C(Y) through
%           the boxed Λ_Y = φ^* Λ^2_X φ_*.
%
% Notation fixed for the batch. The typescript's slashed O is the immersion
% φ : Y → X, and is given as \varphi throughout, with \varphi^{*} and
% \varphi_{*}; the same letter is, in the Hodge-algebra pages, the alternating
% form on V — the two runs are independent and the collision is his. The
% polarisation classes are \xi_{X}, \xi_{Y}; the Lefschetz operators L_{X},
% \Lambda_{X}, L_{Y}, \Lambda_{Y}; the Künneth projectors p^{X}_{i}, p^{Y}_{i}
% (typed as p^X_i, written by hand as p^i_X — both are kept as they stand);
% the primitive parts P^{i}(X). The coproduct sign he draws for a direct sum
% of graded pieces is \coprod on pages 33 to 39 and \oplus where he writes a
% circled plus; both are transcribed as drawn.
%
% Page 39 breaks off on « Mais » and runs into the next batch; page 21 opens
% mid-computation and continues something on a leaf this batch does not hold.
%
% The dating is the inventory's own, for the whole shelfmark.
\documentclass[11pt,a4paper]{article}
% Le préambule commun à toutes les transcriptions du fonds Grothendieck.
%
% Il tient en une idée : **l'incertitude se compose, elle ne se résout pas.**
% Une transcription qui lisse un mot douteux a détruit la seule chose pour
% laquelle on la faisait — un lecteur doit pouvoir savoir, à chaque endroit,
% ce qui était sur le feuillet et ce qui a été ajouté. Les six macros qui
% suivent sont donc l'appareil critique, et non de la décoration.
%
% Les mêmes commandes sont reconnues par `scripts/render.mjs`, qui produit la
% vue de lecture du site. Une macro ajoutée ici doit l'être aussi là-bas,
% faute de quoi le rendu échouera bruyamment — ce qui est voulu : un rendu
% silencieusement faux est pire qu'un rendu absent.

\ProvidesPackage{grothendieck}[2026/08/08 Transcriptions du fonds Grothendieck]

\usepackage{amsmath,amssymb,amsthm}
\usepackage{mathrsfs}
\usepackage{stmaryrd}
% Les diagrammes commutatifs sont partout dans ce fonds : c'est la langue
% même de la géométrie algébrique de Grothendieck, et une transcription qui
% les rendrait en prose ne transcrirait rien.
\usepackage{tikz-cd}
% Français par défaut — c'est la langue des feuillets — mais l'anglais est
% chargé aussi : la traduction et le résumé sont des documents anglais, et la
% typographie française (espaces avant les deux-points) y serait une faute.
% Ces documents basculent par \selectlanguage{english} en tête de corps.
\usepackage[english,french]{babel}
\usepackage{fontspec}
\usepackage{geometry}
\geometry{a4paper,margin=25mm,marginparwidth=28mm}
\usepackage{marginnote}
\usepackage{xcolor}
% ulem, not soul. `soul`'s \st and \ul are built on a character-by-character
% scan that XeLaTeX's font machinery breaks: the document fails with an
% undefined control sequence rather than degrading. ulem does the same two jobs
% and is Unicode-engine safe. `normalem` stops it hijacking \emph, which the
% transcriptions use for Grothendieck's own emphasis.
\usepackage[normalem]{ulem}
\usepackage{hyperref}
\hypersetup{colorlinks=true,linkcolor=black,urlcolor=black,pdfborder={0 0 0}}

% --- Métadonnées du lot -----------------------------------------------------
% Reprises telles quelles par le script de rendu, qui les affiche en tête de
% la vue de lecture. Elles fixent ce que le fichier prétend couvrir : sans
% elles, une transcription flotte sans point d'attache dans un fonds de seize
% mille feuillets.

\newcommand{\folder}[1]{\gdef\grothFolder{#1}}
\newcommand{\batch}[1]{\gdef\grothBatch{#1}}
\newcommand{\leaves}[2]{\gdef\grothFirst{#1}\gdef\grothLast{#2}}
\newcommand{\foldertitle}[1]{\gdef\grothFoldertitle{#1}}
\gdef\grothFolder{?}\gdef\grothBatch{?}\gdef\grothFirst{?}\gdef\grothLast{?}\gdef\grothFoldertitle{}

% --- Le résumé --------------------------------------------------------------
%
% Il ouvre la lecture modernisée, sous le titre « Résumé », et il est composé
% en petit. Ce n'est pas un ornement : c'est le seul passage du document qui
% ne soit pas des mathématiques, et un lecteur doit pouvoir le reconnaître
% avant de l'avoir lu — puis le sauter s'il connaît déjà le sujet, ou s'y
% arrêter s'il ne le connaît pas. Le filet qui le suit marque la reprise du
% corps.

\newenvironment{resume}
  {\small\setlength{\parskip}{0.45em}}
  {\par\medskip\hrule\medskip}

% --- Mots-clés ---------------------------------------------------------------
%
% En anglais, en clôture du résumé de la lecture modernisée : le vocabulaire
% moderne sous lequel chercher ce que les feuillets construisent. Le manifeste
% du site (`npm run manifest`) les extrait pour étiqueter la cote entière —
% cette ligne est la source unique des tags, il n'y a pas de fichier à côté.
\newcommand{\keywords}[1]{\par\smallskip\noindent
  {\footnotesize\textsc{Keywords} --- \textit{#1}}\par}

% --- L'appareil critique ----------------------------------------------------

% Le numéro de feuillet, en marge. C'est l'ancre qui permet de régler un
% désaccord sur une formule en regardant *un* feuillet plutôt qu'une chemise
% de six cents. Le site s'en sert aussi pour tourner les pages du fac-similé
% quand on fait défiler la transcription.
\newcommand{\leaf}[1]{%
  \marginnote{\footnotesize\color{gray!70}\textsf{#1}}%
  \label{leaf:#1}\ignorespaces}

% Illisible : on ne devine pas. Un mot inventé qui se lit comme les autres est
% le pire résultat possible.
\newcommand{\ill}{\textcolor{red!60!black}{[\,\ldots\,]}}

% Lecture douteuse : le mot est donné, et signalé comme incertain.
\newcommand{\uncertain}[1]{\uline{#1}}

% Ajout de l'éditeur — une lettre manquante, un mot sous-entendu.
\newcommand{\add}[1]{\textcolor{blue!50!black}{[#1]}}

% Biffé par Grothendieck lui-même. Ce qu'il a rayé fait partie du document :
% une rature dit souvent où la pensée a changé de direction.
\newcommand{\struck}[1]{\sout{#1}}

% Note du transcripteur, sur l'état matériel du feuillet.
\newcommand{\note}[1]{\par\smallskip{\small\color{gray!60!black}\textit{[#1]}}\par\smallskip}

% Note marginale de Grothendieck, distincte du corps du texte.
\newcommand{\marginal}[1]{\par\smallskip\noindent\hspace{1em}%
  {\small\color{gray!60!black}#1}\par\smallskip}

% --- Titre ------------------------------------------------------------------

\AtBeginDocument{%
  \begin{center}
    {\large\bfseries Fonds Alexandre Grothendieck --- cote n\textsuperscript{o}~\grothFolder}\\[2pt]
    {\itshape\grothFoldertitle}\\[4pt]
    {\small feuillets \grothFirst--\grothLast\ (lot \grothBatch)}\\[2pt]
    {\footnotesize\color{gray!60!black}Transcription établie d'après le fac-similé numérisé
     par l'Université de Montpellier.\\
     Les lectures incertaines sont soulignées, les illisibles notées [\,\ldots\,],
     les ajouts entre crochets.}
  \end{center}
  \vspace{1em}}

\endinput


\folder{16}
\batch{2}
\pages{21}{40}
\dating{1966-1967}
\watermark{Édition de démonstration}
\foldertitle{Formalisme algébrique des correspondances et algèbres de Poincaré-Hodge (cf conjectures standard) : notes manuscrites (s.d.), tapuscrit annoté (s.d.), lettres (1966-1967)}

\begin{document}

\section*{$L^{n}_{0}$, $\Lambda$, et une algèbre de dimension 6 (page 21)}

\note{Feuillet au crayon, écrit vite et biffé aux trois quarts. Les formules
se lisent, la prose qui les relie souvent pas ; l'apparat le dit à chaque
fois. Le titre de la section est de l'éditeur.}

\page{21}
\struck{\ill{}} $\xi^{2n}$ \struck{\ill{}}

\note{Le haut du feuillet est barré de grandes boucles et n'a pas été lu. La
marge supérieure droite porte une indication du même genre, également biffée.}

\[
L^{n}_{0} = L^{n}_{0}\,\pi_{0} = L^{n}_{0}\,\Lambda'^{\,n} L'^{\,n}
 = L^{n}_{0}\,(\beta \Lambda^{n}_{0})(\alpha L^{n}_{0})
\]
\[
= L^{n}_{0}\,\alpha\beta\,\pi_{0} = \alpha\beta\,L^{n}_{0}
\]
donc $\alpha\beta = 1$, \ill{} $\alpha$ et $\beta$ \ill{}

\note{Au-dessus de la première ligne, une note interlinéaire, non lue, qui
paraît poser $L'^{\,n} = \alpha L^{n}_{0}$ et $\Lambda'^{\,n} = \beta \Lambda^{n}_{0}$,
ce que la suite du calcul demande.}

Donc \uncertain{$L'$ de Lefschetz} $\Longrightarrow$ $L^{n}$ \uncertain{un bon}
\ill{} $\xi^{2n}$, \ \ $\overline{\Lambda}$ \uncertain{un bon} \ill{} $\xi^{-2n}$

\note{Les deux mots soulignés sur la page, qu'on lirait volontiers
« une base de », ne se laissent pas assurer ; le reste de la phrase, y compris
ce que $L^{n}$ et $\overline{\Lambda}$ sont censés y être, n'a pas été lu.}

\ill{} \uncertain{On ne demande} \ill{} que $\xi$ \ill{} où \ill{} $\xi$ est
\uncertain{projectif de type fini} \ill{} $k$. Ce \ill{} tel quel, p.~ex.
\ill{} $= \uncertain{\mathcal{L}(M^{*})}$, $M^{*}$ gradué \ill{} degrés
$0 \le i \le 4$ ---

Mais ça \ill{} plus précises, \ill{} de degré 2 \ill{}

\uncertain{Lefschetz}, pas \ill{} de dim.\ 6 \ill{} de degré 2 \ill{} les
relations
\[
L^{4} = 0, \quad L'^{\,2} - LL' = 0, \quad L^{2}L' - L^{3} = 0, \quad L^{3}L' = 0,
\]
donc si \ill{}
\[
1 \ \| \ L,\ L' \ \| \ L^{2},\ LL' \ \| \ L^{3}
\qquad (\deg : \ 0 \quad 2 \quad 4)
\]
$L$ est Lefschetz, bien que $L'$ ne l'est pas, bien que $L'^{\,3} = L^{3}$\,!

\note{Les traits doubles séparent les composantes homogènes, dont les degrés
$0$, $2$, $4$ sont portés dessous ; $L^{3}$, de degré $6$, n'est pas étiqueté
sur la page. « ne l'est pas » est écrit ainsi, sans accord avec « bien que ».}

\section*{Algèbres de Hodge (pages 24 à 28)}

\note{Le titre « Algèbres de Hodge » est de sa main : il figure seul en tête
d'un feuillet par ailleurs vierge (page 30), qui sert de chemise à ces trois
leçons et ne reçoit donc pas de numéro de page ici. Dans tout ce qui suit
$\varphi$, $\psi$, $Q$, $\chi$ sont des formes sur $V$ et sur $W$, et n'ont
rien à voir avec l'immersion $\varphi$ du formulaire qui suit.}

\page{24}
\textbf{Algèbres de Hodge d'ordre 2}

\[
\begin{array}{ccccc}
H^{0} & H^{1} & H^{2} & H^{3} & H^{4} \\
k & V & k\xi + W & \check{V} & k\xi^{2}
\end{array}
\]

\[
\Lambda^{2}V \xrightarrow{\ \alpha\ } k\xi + W, \qquad
\alpha(x \wedge y) = \varphi(x,y)\,\xi + \psi(x,y)
\]
\[
\mathrm{Sym}^{2}(W) \xrightarrow{\ \beta\ } k\xi^{2}, \qquad
\beta(xy) = Q(x,y)\,\xi^{2}
\]

avec $\varphi(x,y)$ forme alternée, $\psi(x,y)$ \uncertain{forme bilinéaire
alternée} à valeurs dans $W$, $Q$ forme bilinéaire symétrique sur $W$.

\struck{\ill{}} La connaissance de $V$, $W$, $\varphi$, $\psi$, $Q$ permet de
\uncertain{reconstituer} $H^{*}$ : \add{à} isom.\ canonique près.


$H^{0} \times H^{i} \to H^{i}$ : clair, \struck{pour} $1 \cdot x = x \cdot 1 = x$.

$H^{1} \times H^{i} \to H^{i+1}$ :
\begin{itemize}
\item \struck{dont} $i = 1$ : $xy = \varphi(x,y)\,\xi + \psi(x,y)$ \quad $x, y \in V$
\item $i = 2$ : $x\xi = \xi x = \tilde\varphi(x) : y \mapsto \varphi(y,x)$, \quad
      $xw = wx = \rho(w,x) : y \mapsto Q(\psi(y,x), w)$ \quad $(x \in V,\ y \in V,\ w \in W)$
\item $i = 3$ : $xy' = -y'x = \struck{\ill{}}\ \langle x, y'\rangle\,\xi^{2}$
      \quad $x \in V$, $y' \in \check{V}$
\item \struck{$i = 4$} \struck{\ill{}}
\end{itemize}

$H^{2} \times H^{i'} \to H^{i'+2}$, $i' = 2$ :
\begin{itemize}
\item $\xi \cdot \xi = \xi^{2}$
\item $\xi w = w\xi = 0$ \quad $(w \in W)$
\item $ww' = Q(w,w')\,\xi^{2}$
\end{itemize}

$(*)$ On pose $\tilde\varphi : V \to \check{V}$, avec
$\langle x, \tilde\varphi(y)\rangle = \varphi(x,y)$.

Quels que soient les choix de $\varphi$, $\psi$, $Q$, \ill{} \struck{donnent}
une loi de composition anticommutative, admettant $1$ comme unité.

\emph{Condition d'associativité} : s'écrit $(xy)(tz) = (tx)(yz)$,
\uncertain{sous la forme}
\[
\varphi(x,y)\varphi(t,z) + Q(\psi(x,y), \psi(t,z))
 = \varphi(t,x)\varphi(y,z) + Q(\psi(t,x), \psi(y,z))
\]

\marginal{exprime aussi que la forme \ill{} $x\,y\,z\,t$ est \ill{}
$(xy)(zt)$ \ill{} $y^{2}$ \ill{} \struck{\ill{}}}

\emph{Condition de} \struck{Lefschetz} :
\begin{itemize}
\item $\xi^{2} : H^{0} \xrightarrow{\ \sim\ } H^{4}$ : tjrs satisfait
\item $\xi : H^{1} \to H^{3}$ : sss $\varphi$ non dégénérée
\end{itemize}

\emph{Condition de Poincaré} :
\begin{itemize}
\item $H^{4} \simeq k$ : tjrs satisfait
\item $H^{0} \times H^{4} \to H^{4}$ et $H^{1} \times H^{3} \to H^{4}$ non dég.\ : idem
\item $H^{2} \times H^{2} \to H^{4}$ non dég.\ : sss $Q$ non dégénérée
\end{itemize}

\note{« tjrs » abrège \emph{toujours}, « sss » \emph{si et seulement si},
« non dég.\ » \emph{non dégénérée}. Les deux critères encadrés sur la page
sont ici mis en valeur par la mise en liste.}

\page{26}
\emph{Cas particulier} : celui où $(xy)(zt) = 0$ \add{quels que soient
$x, y, z, t$}, i.e.\ $(H^{1})^{\struck{\otimes}3} = 0$ dans \uncertain{$H^{4}$}\,!
i.e.\ l'image de $\alpha : \Lambda^{2}V \to k\xi + W$ est \ill{}
\uncertain{isotrope} de $k\xi + W$.

\struck{\ill{}} \uncertain{Par exemple}, si $W$ \ill{}, dans le cas \ill{} des
\ill{} est nécessairement ainsi \ill{}

Plus généralement, considérons la forme 4-linéaire alternée $\chi$ définie par
\[
x\,y\,z\,t = \chi(x,y,z,t)\,\xi^{2} \qquad (x,y,z,t \in V)
\]
Alors la forme \add{bilinéaire} induite par la forme $Q \oplus 1$ (sur
$k\xi + \uncertain{V}$) sur $\Lambda^{2}V$ par $\alpha$, qui définit une
\struck{forme} forme linéaire sur $\mathrm{Sym}^{2}(\Lambda^{2}V)$, se
\uncertain{factorise} \ill{}, $\Lambda^{2}V$ exprimant $\chi$, et \ill{} d'où
les données

\begin{enumerate}
\item[a)] $V$ vect.\ (de dim.\ finie) sur $k$
\item[\struck{b)}] \struck{$\varphi : \Lambda^{2}V \to k$}
\item[b)] $\chi : \Lambda^{4}V \to k$, définissant une forme linéaire
      symétrique $\chi'$ sur $\Lambda^{2}V$
\item[c)] un vectoriel $\tilde W$, avec une forme \struck{quadratique}
      \add{bil.\ sym.} $\tilde Q$, et une application linéaire
      \[
      \Lambda^{2}V \xrightarrow{\ \alpha\ } \tilde W
      \]
      compatible avec $\chi'$, $Q$ \add{i.e.\ disant que $\tilde W$, $\tilde Q$
      \ill{}} \uncertain{déduites} de la forme $\chi'$
\item[d)] un élément $\xi$ de $\tilde W$, tel que $Q(\xi,\xi) = 1$.
\end{enumerate}

\emph{Conditions de non-dégénérescence} :
\begin{itemize}
\item $Q$ non dégénérée
\item $\varphi(x,y) = Q(\alpha(x \wedge y), \xi)$ non dégénérée
\end{itemize}

\page{28}
\struck{Exemple. Prenons pour \ill{} dég.\ \ill{}}
\begin{itemize}
\item $\varphi : \Lambda^{2}V \to k$ non dégénérée, i.e.\ $\varphi \in \Lambda^{2}\check{V}$
\item \struck{$\check\varphi \in \Lambda^{2}\check{V}$ \ill{}} \quad
      $\chi = \varphi^{2} \in \Lambda^{4}\check{V}$, \quad $\chi : \Lambda^{4}V \to k$
\end{itemize}
\struck{Alors $\chi$ \ill{} forme \ill{} sur $\Lambda^{2}V$ \ill{}}

Ex. Prenons $\chi'$ non dégénérée, \struck{\ill{}} $\alpha = \mathrm{id}$,
i.e.\ $\tilde W = \Lambda^{2}V$. Alors on choisit un $\xi \in \tilde W = \Lambda^{2}V$
donné n'importe quelle $\varphi$ voulue, à un facteur scalaire, donc n'importe
quelle \uncertain{$\check\varphi \in \Lambda^{2}\check{V}$} non dégénérée, à un
facteur scalaire. Considérons alors $\check\varphi$ comme élément de
\struck{$\Lambda^{2}\check{V}$} l'algèbre $\check\Lambda V \otimes \check\Lambda V$
\add{en écrivant $\Lambda^{2}V \subset \uncertain{V} \otimes V$} et
$\check\varphi^{\,2}$ comme élément de
$\Lambda^{2}\check{V} \otimes \Lambda^{2}\check{V}$ (qui $\in \Lambda^{4}\check{V}$)
$\subset \check\Lambda V \otimes \check\Lambda V$ ; le carré de $\check\varphi$ de
$\Lambda^{2}V$ (à l'image de $\check\varphi^{\,(2)}$ \uncertain{prise} par
$\Lambda^{2}V \otimes \check\Lambda V \to \Lambda^{4}V$). Utilisant la forme
$Q = Q_{\chi'}$, définie par $\chi'$, on voit que $\check\varphi^{\,(2)}$
définit un endom.\ de $\Lambda^{2}V$, soit $u$. Je dis que l'orbite sous $u$
des éléments $\xi$ et $\check\varphi$ de $\Lambda^{2}V$ \uncertain{peuvent}
\struck{jamais} être infinis, si $\dim V \ge 4$ (ou 6\,?) ($W$ \ill{} de dim.\
paire) \struck{\ill{}} et \uncertain{plusieurs} \ill{} de dim.\ non bornés
(qui dès $V$ \ill{}).

\note{La lettre majuscule que porte l'inclusion $\Lambda^{2}V \subset \ldots
\otimes V$ se lit aussi bien $N$ que $V$ ; on a écrit $V \otimes V$, qui est
ce que l'énoncé demande, et on le signale. Le passage final, sur la finitude
de l'orbite, est une esquisse rapide dont la syntaxe ne se reconstitue pas.}

\section*{Formulaire de $L$, $\Lambda$ (pages 32 à 39)}

\note{Le titre est de sa main, porté en tête d'un feuillet vierge (page 31)
qui ouvre cette suite et ne reçoit pas de numéro ici ; le tapuscrit de la page
32 porte le même titre. Le tapuscrit écrit l'immersion $Y \to X$ avec un O
barré ; on l'a rendue partout par $\varphi$, comme le fait sa propre main dans
la suite manuscrite. Les corrections portées à l'encre sur le tapuscrit sont
en \add{} ou en \struck{}, ses renvois en \marginal{}.}

\page{32}
\subsection*{Formulaire de $L$, $\Lambda$.}

$X$ variété projective lisse connexe polarisée par $\xi_{X} \in H^{2}(X)$,
$L_{X}$ opération de cup-produit associée, $\Lambda_{X}$ opération de degré
$-2$ correspondante, $p^{X}_{i}$ les projections de $H^{*}(X)$ sur les
morceaux primitifs $P^{i}(X)$ pour $0 \le i \le 2n$, $Y$ une section
hyperplane lisse de $X$, $\xi_{Y} \in H^{2}(Y)$ la self-intersection (induite
par $\xi_{X}$), d'où $L_{Y}$, $\Lambda_{Y}$, $p^{i}_{Y}$ ;
$\varphi : Y \to X$ l'immersion, d'où $\varphi^{*}$ et $\varphi_{*}$.

A) \quad On a
\begin{align}
\mathrm{Im}\ L_{X} &= \bigcap_{0 \le i \le n} \mathrm{Ker}\ p^{X}_{i}, &
\mathrm{Ker}\ L_{X} &= \sum_{n \le i \le 2n} \mathrm{Im}\ p^{X}_{i}, \\
\mathrm{Im}\ \Lambda_{X} &= \bigcap_{n \le i \le 2n} \mathrm{Ker}\ p^{X}_{i}, &
\mathrm{Ker}\ \Lambda_{X} &= \sum_{0 \le i \le n} \mathrm{Im}\ p^{X}_{i},
\end{align}

d'où en particulier les formules
\begin{align}
p^{X}_{i} L_{X} &= 0 \ \text{ si } 0 \le i \le n, &
L_{X} p^{X}_{i} &= 0 \ \text{ si } n \le i \le 2n, \\
\uncertain{p^{X}_{i}}\, \Lambda_{X} &= 0 \ \text{ si } n \le i \le 2n, &
\Lambda_{X} p^{X}_{i} &= 0 \ \text{ si } 0 \le i \le n,
\end{align}

et en particulier
\begin{equation}
p^{X}_{n} L_{X} = L_{X} p^{X}_{n} = p^{X}_{n} \Lambda_{X} = \Lambda_{X} p^{X}_{n} = 0 .
\end{equation}

\note{Le $p^{X}_{i}$ de la formule 2') est surchargé à l'encre ; la correction
n'a pas été lue et l'indice typographié est conservé.}

B) \quad On a
\begin{equation}
\Lambda_{X} L_{X} = \mathrm{id}_{X} - \sum_{n \le i \le 2n} p^{X}_{i}, \qquad
L_{X} \Lambda_{X} = \mathrm{id}_{X} - \sum_{0 \le i \le n} p^{X}_{i},
\end{equation}

\note{Les deux bornes de sommation de 4) sont repassées à l'encre sur le
tapuscrit ; ce sont les valeurs corrigées qui sont transcrites.}

\struck{d'où compte tenu de 1') et 2') les formules}

\struck{4') $\Lambda_{X} L^{2}_{X} = \bigl(\mathrm{id}_{X} - \sum_{0 \le i \le 2n} p^{X}_{i}\bigr) L_{X}$, \quad
$L^{2}_{X} \Lambda_{X} = L_{X}\bigl(\mathrm{id}_{X} - \sum_{0 \le i \le 2n} p^{X}_{i}\bigr)$}

\struck{4'') $L_{X} \Lambda_{X} L_{X} = L_{X}\,\ill{}$}

\marginal{et formules « transposées ».}

C) \quad On a
\begin{equation}
\varphi_{*} \Lambda_{Y} \varphi^{*} = \mathrm{id}_{X} - \sum_{0 \le i \le 2n} p^{X}_{i}
\end{equation}

\marginal{par 4) : $= (\Lambda_{X} L_{X} + L_{X} \Lambda_{X}) - \mathrm{id}_{X} + p^{X}_{n}$}

d'où, \struck{compte tenu de 4''), 1') et 2')} \add{posant}
\begin{equation}
T_{X} = \mathrm{id}_{X} + \varphi_{*} \Lambda_{Y} \varphi^{*},
\end{equation}

\struck{5') $L_{X} \Lambda_{X} L_{X} = \ill{}$}

\marginal{(donc $\Lambda_{X} L_{X} + L_{X} \Lambda_{X} = T_{X} - p^{X}_{n}$),}

d'où compte tenu de \struck{3}\add{*}) la formule
\begin{equation}
L^{2}_{X} \Lambda_{X} + 2\,L_{X} \Lambda_{X} L_{X} + \Lambda_{X} L^{2}_{X}
 = \struck{L_{X}(\mathrm{id}_{X} + \varphi_{*}\Lambda_{Y}\varphi^{*}) + (\mathrm{id}_{X} + \varphi_{*}\Lambda_{Y}\varphi^{*})L_{X}}
 \ \add{L_{X} T_{X} + T_{X} L_{X}},
\end{equation}
avec $T_{X}$ \struck{$= \mathrm{id}_{X} + \varphi_{*}\Lambda_{Y}\varphi^{*}$}
\add{donné par (6)}.

\textbf{Corollaire.} $(A(X \times X)$ et $C(Y))$ implique $C(X)$.

\marginal{On y peut remplacer $A(X \times X)$ par $A^{0}(X \times X)$
(correspondances « en dim.\ critique »
$H^{\uncertain{(2n+1)}-2} \to H^{\uncertain{(2n)}+2}$)}

\note{Les exposants de cette dernière note marginale sont surchargés et ne se
laissent pas assurer ; ils sont donnés tels qu'ils se lisent, sous réserve.}

\page{33}
\uncertain{Prop.}
\begin{itemize}
\item $P^{i}(X) \xrightarrow{\ \sim\ } P^{i}(Y)$ \quad si $i \le n-2$
\item $P^{n-1}(X) \hookrightarrow P^{n-1}(Y)$ \quad \struck{si $i = n-1$}
\item $P^{i}(X) = \mathrm{Ker}\ \varphi^{i}$ \quad si $i \ge n$
\end{itemize}

\begin{tikzcd}
\struck{H}^{i}(X) \arrow[r, "L_{X}^{n-i+1}"] \arrow[d, "\varphi^{i}"'] & H^{2n-i+2}(X) \\
H^{i}(Y) \arrow[r, "L_{Y}^{n-i}"] & \uncertain{H^{2(n-1)-i}}(Y) \arrow[u, "\varphi_{*}"']
\end{tikzcd}

\note{Les deux flèches verticales portent chacune un $\simeq$ ; l'exposant de
$\varphi_{*}$ et celui du $H$ en bas à droite sont surchargés et douteux. Le
crochet de droite porte la condition $i \le n-2$.}

\[
i \le n-2 \ \Longrightarrow \ 2(n-1) - i \ge n
\]

\begin{tikzcd}
H^{j}(X) \arrow[d, "\varphi^{*}"'] \arrow[rd, "L_{X}"] & \\
H^{j}(Y) \arrow[r, "\sim"] & H^{j+2}(X)
\end{tikzcd}
\qquad $j \ge n$

\struck{\ill{}} \marginal{Non $\to$}
\[
H^{d}(X) = \struck{\ill{}} \coprod_{\substack{i+j \le n \\ i,j \ge 0}} L^{j}_{X} P^{i}(X)
 \ \oplus \coprod_{j \ge n} P^{j}(X)
\]
\[
H^{d}(Y) = \coprod_{\substack{i+j \le n-1 \\ i,j \ge 0}} L^{j}_{Y} P^{i}(Y)
\]

Car $\varphi^{*}$ \uncertain{induit} un \ill{} d'\uncertain{objets bigradués}
\[
\coprod_{\substack{i+j \le n-1 \\ i,j \ge 0}} L^{j}_{X} P^{i}(X)
\]
\[
P^{i}(X) \longrightarrow H^{d}(Y) = \coprod_{\substack{i+j \le n-1 \\ i,j \ge 0}} L^{j}_{Y} P^{i}(Y) \ ;
\]

\struck{\ill{}} : \ill{}, un \uncertain{bihomogène}, i.e.
\[
L^{0}_{X} P^{n-1}(X) = P^{n-1}(X) \longrightarrow \struck{\ill{}}\ \varphi^{*} L^{0}_{Y} P^{n-1}(Y) = P^{n-1}(Y).
\]
cf.\ \uncertain{2} : \struck{celui de $L_{X}$, i.e.} \struck{$\coprod_{j \ge n} P^{j}(X)$} \struck{\ill{}}

Car le noyau de $\varphi^{*}$ \ill{} \struck{\ill{}}. Le \uncertain{conoyau} \ill{}
\[
E^{n-1} = H^{n-1}(Y)/\mathrm{Im}\ H^{n-1}(X) = P^{n-1}(Y)/\mathrm{Im}\ P^{n-1}(X)
\]
\[
\simeq \mathrm{Ker}\bigl(\varphi_{*} : H^{n-1}(Y) \to H^{n+1}(X)\bigr)
\]

\page{35}
\uncertain{On aura donc}, \ill{} :
\[
\text{(b)} \qquad H^{d}(Y) \simeq
 \coprod_{\substack{i+j \le n-1 \\ i,j \ge 0}} L^{j}_{X} P^{i}(X) \ \oplus\ E^{n-1}
\]

L'opération
\[
\varphi_{*} : H^{d}(Y) \longrightarrow H^{d+2}(X)
\]
s'interprète donc : $L_{X}$ sur le \uncertain{morceau}
$\coprod_{i+j \le n-1,\ i,j \ge 0} L^{j}_{X} P^{i}(X)$ des premiers morceaux de
$\alpha$), (et est \struck{\ill{}} \ill{}), et il est \emph{nul} sur le
morceau \struck{$\coprod_{\ill{}} L^{j}_{X} P^{i}(X) \oplus E^{n-1}$},
\struck{qui correspond} \struck{$\sum_{n-1 \le i \le 2(n-1)} P^{i}(Y)$}. Ainsi

\textbf{Corollaire.}
\[
\mathrm{Ker}\ \varphi_{*} \ \struck{=\ \ill{}}\ \simeq\ \mathrm{Ker}\ E^{n-1}(Y)
\]
\[
\mathrm{Im}\ \varphi_{*} = \mathrm{Im}\ \varphi_{*}\varphi^{*} = \mathrm{Im}\ L_{X}
 = \bigcap_{0 \le i \le n} \mathrm{Ker}\ p^{X}_{i}
\]
car $H^{d}(Y) = \mathrm{Im}\ \varphi^{*} + E^{n-1}$ (la dernière formule, car
\ill{} $E^{n-1}$, \ill{} $\varphi_{*}$ s'annule sur $E^{n-1}$ \ill{} ---
$\varphi_{*}\varphi^{*} = L_{X}$).

\marginal{\struck{\ill{}} \uncertain{variantes}}

On a
\[
\varphi^{*} L_{X} = L_{Y} \varphi^{*}, \qquad L_{X} \varphi_{*} = \varphi_{*} L_{Y}
\]
\[
p^{X}_{i} \varphi_{*} = 0 \ \text{ si } 0 \le i \le n, \qquad
p^{X}_{j} \varphi_{*} = \varphi_{*} p^{Y}_{j-2} \ \text{ si } \struck{n+1} \le i \le 2n
\]
\[
\varphi_{*} \varphi^{*} = L_{X}, \qquad \varphi^{*} \varphi_{*} = L_{Y}
\]
formules de projection.

Pour la dernière formule, \ill{} \ill{}, donc il suffit de vérifier
$(L_{X}\varphi_{*})\varphi^{*} = (\varphi_{*}L_{Y})\varphi^{*}$, i.e.\
$L_{X}(\varphi_{*}\varphi^{*}) = \varphi_{*}(L_{Y}\varphi^{*})$ \uncertain{car
les deux membres sont} $L^{2}_{X}$.

\page{37}
A) \quad Moyennant l'\uncertain{identification} \ill{} ($\alpha$), l'opération
$L_{Y}$ \uncertain{s'interprète} : $L_{X}$ sur le morceau
$\coprod_{i+j \le n-2,\ i,j \ge 0} L^{j}_{X} P^{i}(X)$, et \ill{} le morceau
\[
\coprod_{\substack{i+j = n-1 \\ i,j \ge 0}} L^{j}_{X} P^{j}(X) \ \oplus\ E^{(n-1)}
 \ = \ \coprod_{i \ge n+1} P^{j}(X) \ \oplus\ E^{(n-1)}
\]
(qui n'a d'\ill{} \ill{} \uncertain{isomorphe} par
$\coprod_{\ill{} \le i \le 2(n-1)} P^{i}(Y)$, \ill{} il y a droit).

L'opération $\Lambda_{Y}$ \uncertain{est} nulle sur $E^{n-1}$, et coïncide sur
le \uncertain{morceau} $\coprod_{i+j = n-1,\ i,j \ge 0} L^{j}_{X} P^{i}(X)$
\ill{} avec $\Lambda_{X}$. Car \uncertain{vectoriellement}
\[
\coprod_{0 \le i \le n-1} P^{i}(X) + E^{n-1} \ \simeq\
\coprod_{0 \le i \le n-1} P^{i}(Y),
\]
\ill{} et \ill{} $\Lambda_{Y}$ \ill{}
\[
\longrightarrow \coprod_{\substack{i+j \le n-2 \\ i,j \ge 0}} L^{j-1}_{X} P^{i}(X) .
\]

\begin{itemize}
\item \struck{\ill{}} $\varphi^{*} \Lambda_{X} = \Lambda_{Y} \varphi^{*} + \ill{}$
\item $\varphi_{*} \Lambda_{Y} = \Lambda_{X} \varphi_{*}$
\item Par la \uncertain{même vérification}, i.e.\ sur les
      $L^{n-i}_{X} P^{i}(X)$ pour $i \le n-1$ ; $P^{j}(X)$ pour $j \ge n$,
      i.e.\ \ill{} ; $\varphi^{*}$ \ill{}, les \ill{} \ill{}
\end{itemize}

\marginal{$\Lambda^{m}_{Y}\varphi^{*} = \varphi^{*}\Lambda^{m}_{X}(\mathrm{id} - \sum \ill{})$ \ill{}}

\[
\Lambda_{Y} \varphi^{*} = \varphi^{*} \Lambda_{X}\Bigl(\mathrm{id}_{X} - \sum_{n+1 \le j \le 2n} p^{j}_{X}\Bigr)
 = \Bigl(\mathrm{id}_{Y} - \struck{\sum_{n+1 \le j \le 2(n-1)}} p^{j}_{X}\Bigr) \varphi^{*} \Lambda_{X}
\]
\[
\varphi^{*} \Lambda_{X} p^{j}_{X} = p^{j-2}_{Y} \varphi^{*} \Lambda_{X} \struck{\ill{}}
\quad \text{si } j \ge n+1
\]
\[
\varphi_{*} \Lambda_{Y} = \Bigl(\mathrm{id}_{X} - \sum_{0 \le i \le n-1} p^{i}_{X}\Bigr) \Lambda_{X} \varphi_{*}
 = \Lambda_{X} \varphi_{*} \Bigl(\mathrm{id}_{Y} - \sum_{0 \le i \le n-1} p^{i}_{Y}\Bigr)
\]
\[
p^{i}_{X} \Lambda_{X} \varphi^{*} = \Lambda_{X} \varphi^{*} p^{i}_{Y} = p^{i}_{X} \varphi_{*} \Lambda_{Y}
\quad \text{si } 0 \le i \le n-1
\]

\page{39}
On en \uncertain{conclut} des \uncertain{premières} \ill{} que
\[
\varphi_{*} \Lambda_{Y} \varphi^{*}
 = \underbrace{\varphi_{*}\varphi^{*} \Lambda_{X}}_{L_{X}\Lambda_{X}}
   \Bigl(\mathrm{id}_{X} - \sum_{n+1 \le j \le 2n} p^{j}_{X}\Bigr)
 = \underbrace{\mathrm{id}_{X} - \sum_{0 \le j \le 2n} p^{j}_{X}}_{q}
\]

\note{Sous l'accolade de gauche, $L_{X}\Lambda_{X}$, et sous elle encore
$= \mathrm{id}_{X} - \sum_{0 \le i \le 2n} p^{X}_{i}$ ; l'accolade de droite
nomme cet opérateur $q$.}

\add{Plus gén.}
\[
\varphi_{*} \Lambda^{m}_{Y} \varphi^{*}
 = \struck{\ill{}} \Bigl(\mathrm{id} - \sum_{0 \le i \le n} p^{i}_{X}\Bigr)
   \Lambda^{m-1}_{X} \Bigl(\mathrm{id} - \sum_{n+1 \le j \le 2n} p^{j}_{X}\Bigr)
 = q\,\Lambda^{m-1}_{X}\,q
\]
où $q = \mathrm{id} - \sum_{0 \le i \le 2n} p^{i}_{X}$.

Ainsi :
\[
\varphi^{*} \Lambda_{X} \varphi_{*} = \mathrm{id}_{Y} - p_{E}
\]
où $p_{E}$ est \uncertain{la projection sur} la partie
\uncertain{supplémentaire} $E^{n-1}$, d'\ill{}, \ill{} multiplié par
$\Lambda_{Y}$ à \uncertain{gauche}, comme $\Lambda_{Y} p_{E} = 0$ :
\[
\Lambda_{Y} = \Lambda_{Y} \varphi^{*} \Lambda_{X} \varphi_{*}
 = \varphi^{*} \Lambda_{X} \Bigl(\mathrm{id}_{X} - \sum_{n+1 \le j \le 2n} p^{j}_{X}\Bigr) \Lambda_{X} \varphi_{*}
\]
donc
\[
\boxed{\ \Lambda_{Y} = \varphi^{*} \Lambda^{2}_{X} \varphi_{*}\ }
\]

\note{L'accolade au-dessus du facteur central porte $= 0$ pour le terme
soustrait ; c'est ce qui donne la formule encadrée.}

Donc $C(X)$ \struck{\uncertain{implique}} $C(Y)$.

Soient $p^{i,j}_{X}$ les \uncertain{projections} \ill{} sur les
$L^{j}_{X} P^{i}(X)$ ($i+j \le n$, $i,j \ge 0$). Alors
\[
\varphi_{*}\, p^{i,j}_{Y}\, \varphi^{*} = p^{i,j+1}_{X}
\qquad \text{pour } 0 \le i+j \le n-1,\ i,j \ge 0 \ \struck{\ill{}}
\]
dans le \uncertain{cadre} \uncertain{de} $D(Y)$, ce qui implique que les
$p^{i,j}_{X}$ sont \ill{} algébriques pour \ill{} de $i+j \le n$, $i \ge 0$ et
$j \ge 1$, i.e.\ les mêmes \ill{} \ill{} $p^{i}_{X}$ pour
$0 \le i \le \struck{2n}$. Mais

\note{Le feuillet s'arrête sur ce « Mais ». La page 40 étant un tapuscrit
étranger, la suite est la page 41, qui reprend la même discussion de $C(X)$ et
de $C(Y)$.}

\end{document}
